%%%%%%%%%%%%%%%%
% This is an example CV created using altacv.cls (v1.4, 12 Apr 2021) written by
% LianTze Lim (liantze@gmail.com), based on the
% Cv created by BusinessInsider at http://www.businessinsider.my/a-sample-resume-for-marissa-mayer-2016-7/?r=US&IR=T
%
%% It may be distributed and/or modified under the
%% conditions of the LaTeX Project Public License, either version 1.3
%% of this license or (at your option) any later version.
%% The latest version of this license is in
%%    http://www.latex-project.org/lppl.txt
%% and version 1.3 or later is part of all distributions of LaTeX
%% version 2003/12/01 or later.
%%%%%%%%%%%%%%%%

%% Use the "normalphoto" option if you want a normal photo instead of cropped to a circle
% \documentclass[10pt,a4paper,normalphoto]{altacv}

\documentclass[10pt,a4paper,ragged2e,withhyper]{altacv}

%% AltaCV uses the fontawesome5 package.
%% See http://texdoc.net/pkg/fontawesome5 for full list of symbols.

% Change the page layout if you need to
\geometry{left=1cm,right=1cm,top=1cm,bottom=1cm,columnsep=0.5cm}

% The paracol package lets you typeset columns of text in parallel
\usepackage{paracol}


% Change the font if you want to, depending on whether
% you're using pdflatex or xelatex/lualatex
\iftutex
  % If using xelatex or lualatex:
  \setmainfont{Lato}
\else
  % If using pdflatex:
  \usepackage[default]{lato}
\fi

\newcommand{\scholargoogle}[1]{\href{#1}{Google Scholar}}
\NewInfoField*{lattes}{\faGraduationCap}

% Change the colours if you want to
\definecolor{heading_color}{HTML}{5E5C4B}
\definecolor{heading_rule}{HTML}{5E5C4B}
\definecolor{emphasis_color}{HTML}{012E0D }
\definecolor{body_color}{HTML}{020F05}
\definecolor{subheading_color}{HTML}{0B1268}
\colorlet{name}{black}
\colorlet{tagline}{black}
\colorlet{heading}{heading_color}
\colorlet{headingrule}{heading_rule}
\colorlet{subheading}{subheading_color}
\colorlet{accent}{heading_color}
\colorlet{emphasis}{emphasis_color}
\colorlet{body}{body_color}

% Change the bullets for itemize and rating marker
% for \cvskill if you want to
\renewcommand{\itemmarker}{{\small\textbullet}}
\renewcommand{\ratingmarker}{\faCircle}

%% Use (and optionally edit if necessary) this .cfg if you
%% want to use an author-year reference style like APA(6)
%% for your publication list
\input{pubs-authoryear.cfg}

%% Use (and optionally edit if necessary) this .cfg if you
%% want an originally numerical reference style like IEEE
%% for your publication list
% \input{pubs-num}

%% Publicações mantidas neste arquivo.
\addbibresource{publications.bib}

\begin{document}
\name{Hugo Abreu}
\tagline{Desenvolvedor Sênior (Full Stack/Embedded) | Especialista em IA | MSc em Sistemas}
%% You can add multiple photos on the left or right
%%\photoR{2.5cm}{phtoo}
\personalinfo{%
  % Not all of these are required!
  % You can add your own with \printinfo{symbol}{detail}
  \email{hugoabreu1002@gmail.com}
   \phone{(81)97915-3699}
% \mailaddress{Rua Dr. Enéas de Lucena}
  \location{Rua Wagner 1340, Manaus - AM}
  \linkedin{https://www.linkedin.com/in/hugoabreumendes/}
  \scholargoogle{https://scholar.google.com/citations?user=CUtwvawAAAAJ&hl=en}
  \orcid{0000-0002-0474-2680}
  \lattes{Lattes}{http://lattes.cnpq.br/0341866422719760}
}

\makecvheader

%% Depending on your tastes, you may want to make fonts of itemize environments slightly smaller
\AtBeginEnvironment{itemize}{\small}

%% Set the left/right column width ratio to 6:4.

% Start a 2-column paracol. Both the left and right columns will automatically
% break across pages if things get too long.
\begin{paracol}{2}

\cvsection{Experiência Profissional}

\cvevent{Desenvolvedor de Software Sênior}{SIDIA}{Jun. 2022 -- Atual}{Manaus, AM}
\begin{itemize}

    \item Desenvolvimento de PoCs como: Estimativa de Posição de TVs em ambientes internos para prevenção de pirataria e sugestão de conteúdo, análise de comportamento de usuários para VoD via enxame de agências e recuperação de conteúdo. Modelagem de redes de grafos para navegação de TV para deficientes visuais via controle remoto.

    \item Python/NodeJS/React: Animações de IA Generativa baseadas em movimentos humanos + FastAPI \& React. Estimativa de pose e detecção de objetos com mediapipe e yolo, respectivamente. Processamento de vídeo e detecção de similaridade. NLP (Spacy, NLTK) para geração de dataset PT-BR x LIBRAS. Uso de modelos LSTM (TensorflowJS) e T5 Transformers (HuggingFace) para tradução Seq2Seq. Transcrição de áudio com Whisper e Vosk e libs nativas em TizenOS. Também criando uma lib completa em C\# do zero para tradução pt-BR para LIBRAS baseada em mudanças de estrutura semântica e sintática de ambas as línguas.

    \item Trabalho na implantação da solução de tradução de português para LIBRAS para transmissão DTV+ (TV 3.0).
    
    \item Desenvolvimento de um app NodeJS + React + Electron (front-end). Uso de Python para tarefas de descrição de imagem e vídeo com modelos CLIP e BLIP, e aplicação de modelos NLP e Transformers para paráfrase (HuggingFace).
    
    \item C++ (20, 17, 11), C\# (.NET 8): Desenvolvimento DTVPlus para TVs Tizen. Manutenção, depuração e adição de novos recursos a bibliotecas e aplicações C/C++ e C\# para Tizen, especificamente para câmeras e dispositivos IoT.
    
    \item Automação de testes de desempenho de firmware com um app C++/Bash para Linux embarcado (TizenOS), usando Python Flask para APIs e renderização front-end, Docker + PostgreSQL para o banco de dados do app.
    
    \item Aplicação de deep learning para verificar imagens de checkpoint e classificar automaticamente KPIs de desempenho. Uso de Python Django RestAPI para coleta de dados, OpenCV e Scikit-learn + TensorFlow para gerar classificações.
    
    \item Escrita de artigos científicos e patentes.
    
\end{itemize}
\divider

\cvevent{Prompt Engineer}{Outlier}{Jun. 2023 -- Fev. 2025}{Remoto | Freelance}

Aprimorando Grandes Modelos de Linguagem para necessidades específicas de treinamento:

\begin{itemize}
        
    \item Revisão de código gerado por LLMs em múltiplas linguagens de programação (Go, Python, C, C++, C\#, Java, JavaScript, TypeScript) para atender a prompts específicos.

    \item Especialização em ciência de dados e treinamento de grandes modelos de linguagem para resolver diversos desafios relacionados a prompts.

    \item Trabalho em projetos de aprendizado por reforço com feedback humano relacionados à geração de prompts, geração de código, revisão de código e correção de código para grandes modelos de linguagem.

    \item Aplicação de técnicas adversárias em chatbots para detectar possibilidade de geração de malware.

    \item Detecção de malware gerado por LLMs através de revisões de código e análise de intenção de prompts.

    \item Geração de novos prompts do zero para testar as capacidades do LLM e revisão de prompts gerados por outros para ajustar em termos de usabilidade e níveis de dificuldade.

    \item Revisão de rotulagem gerada por outros como prompt engineer de confiança.
    
\end{itemize}
\divider

\cvevent{Engenheiro de Sistemas}{FITec}{Fev. 2019 -- Maio. 2022}{Recife, PE}
\begin{itemize}

    \item C\# .Net para sistemas embarcados (WindowsCE) e Windows Apps para comunicação e integração de sistemas, usados para supervisão e ajuste de rotas de trens.

    \item Desenvolvimento de ferramentas de controle de qualidade e automação para fabricantes de computadores, notebooks e all-in-ones. BASH e Python: FastAPI.
    
    \item Implantação e atualização de servidores Nokia VitalQIP IPAM (IP Address Management) e DNS em servidores RHEL.
    
    \item Automação da importação massiva de dados de rede no VitalQIP, mantendo a estrutura de rede e inserindo informações de clientes finais.
    
    \item Desenvolvimento de soluções de automação em Python, Java e BASH, para integração em servidores RHEL, comunicando-se com roteadores Nokia 77XX.
    
    \item Desenvolvimento de scripts em Apache Velocity na ferramenta Nokia NSP para automatizar a configuração de roteadores Nokia 77XX.
    
    \item Implantação do Nokia NSP para gerenciamento de roteadores em servidores RHEL.
    
    \item Implementação de Nokia AAA (protocolo Radius) em servidor CentOS / RHEL.
    
    \item Análise e correção de vulnerabilidades de segurança em sistemas RHEL, incluindo hardening de servidores.
    
    \item Fornecimento de soluções para redes IP, incluindo upgrades lógicos na infraestrutura de rede e garantia da consistência de pools IP e topologia de rede.
    
    \item Personalização de distribuição Linux baseada em Debian para instalação OEM.
    
    \item Resolução de problemas multidisciplinares utilizando aprendizado de máquina, otimização, estatística e outras ferramentas matemáticas e computacionais, como processamento digital de sinais e imagens, com bibliotecas Python como NumPy, Scikit-learn, SciPy e Pandas.
    
    \item Desenvolvimento de aplicações web (API) e microsserviços com modelos computacionais usando SpringBoot e Flask. Uso da ferramenta de teste Selenium para Web Scraping.
    
    \item Implementação de caderno de teste na etapa de desenvolvimento de um medidor inteligente (Protocolo DLMS). Automação da estrutura de objetos COSEM em XML usando Python.
    
\end{itemize}
\divider

\cvevent{Monitor de Aulas de CLP}{CTG UFPE}{Ago. 2018 --  Dez. 2018}{Recife, PE}
\begin{itemize}

    \item Assistência em aulas de CLP para os cursos de Engenharia Elétrica e Controle \& Automação na UFPE.
    \item Preparação de aulas e resolução de exercícios envolvendo o uso e programação de CLPs SIEMENS S7-1200.
    \item Instrumentação industrial usando motores de passo, motores AC \& DC, encoders ópticos e inversores de frequência.
    
\end{itemize}

\divider

\cvevent{Estágio de Pesquisa}{FACEPE}{Ago. 2017 -- Jul. 2018}{Recife, PE}
\begin{itemize}

    \item Desenvolvimento de equipamento de imagem por infravermelho próximo para diagnóstico de lesões de pele.
    \item Envolvimento em prototipagem mecânica usando SolidWorks, projeto de circuitos fotoeletrônicos e condução de processamento de imagem e sinal com Python e MATLAB.
    \item Processamento digital de espectros de absorbância de luz infravermelha da pele, utilizando PCA e SVM para diagnóstico de doenças autoimunes raras.
    \item Desenvolvimento de um dispositivo de imagem usando LEDs infravermelhos e fotossensores. Condução de processamento digital de imagem de dentes bovinos para estudar radiação e espalhamento infravermelho, e avaliar qualidade de imagem.
    \item Desenvolvimento de um sistema supervisório baseado em LabView para caracterização de amostras aquosas, comunicando via OPC com MATLAB e hardware de metrologia da National Instruments, MyDAQ.
    
\end{itemize}

\divider

\cvevent{Estágio}{WT Engenharia}{Out. 2017 -- Maio. 2018}{Recife, PE}
\begin{itemize}

    \item Manutenção e montagem de módulo eletrônico embarcado de comunicação RF para detecção de orvalho, como protocolo seguro, para trabalhar com linhas de transmissão de alta tensão energizadas.
    \item Projeto IoT aplicado à atuação de chaves seccionadoras usando microcontroladores (Arduino) e rede WIFI (ESP8266), protocolo mestre-escravo TCP/IP.
    
\end{itemize}


\switchcolumn

\cvsection{Publicações Científicas}

\nocite{*}

% \printbibliography[heading=pubtype,title={\printinfo{\faBook}{Books}},type=book]

\divider

\printbibliography[heading=pubtype,title={\printinfo{\faFile*[regular]}{Artigos Científicos em Periódicos}}, type=article]

\divider

\printbibliography[heading=pubtype,title={\printinfo{\faUsers}{Anais de Conferências}},type=inproceedings]

\divider

\printbibliography[heading=pubtype,title={\printinfo{\faUsers}{Patentes}},type=misc]

%% Switch to the right column. This will now automatically move to the second
%% page if the content is too long.

\cvsection{Educação}

\cvevent{Mestrado em Engenharia de Sistemas}{Universidade de Pernambuco}{Mar. 2020 -- Jan. 2022}{Recife, BR}

\divider

\cvevent{Especialização em Inteligência Artificial}{UPE}{Fev. -- Nov. 2019}{Recife, BR}

\divider

\cvevent{Engenharia de Controle e Automação}{Universidade Federal de Pernambuco}{2013 -- 2018}{Recife, BR}

\divider

\cvevent{Engenharia Eletrônica}{ENSEEIHT}{2015 -- 2016}{Toulouse, FR}

\switchcolumn

\cvsection{Programação \& Tecnologia}

\cvtag{Python}
\cvtag{Pytorch}
\cvtag{TensorFlow}
\cvtag{LLM / Transformers}
\cvtag{SciKitLearn}
\cvtag{Pandas}
\cvtag{Visão Computacional}
\cvtag{Séries Temporais}
\cvtag{Prompt Engineering}

\cvtag{C/C++ \& C\#}
\cvtag{STL}
\cvtag{Boost}
\cvtag{Glib}
\cvtag{GStreamer}

\cvtag{Linux \& Scripts Bash}
\cvtag{JavaScript}
\cvtag{React}
\cvtag{JAVA}

\cvsection{Certificações}

\cvtag{Docker: Criando e Gerenciando Containers}
\cvtag{Testes em .NET: Criando Testes de Unidade com xUnit}
\cvtag{Developing on AWS (3.4.0)}
\cvtag{.NET: Criando uma API Web com ASP.NET Core}
\cvtag{Hugging Face: Aprimorando Eficiência de Modelos com Optimum}

\cvsection{Línguas Estrangeiras}

\cvskill{Português (nativo)}{5}

\cvskill{Inglês}{5}

\cvskill{Francês}{4}

\cvskill{Espanhol}{2}



\end{paracol}
\end{document}